\TNchapter[By the end of Part III you can read a vendor's score. By the end of this chapter you can decide what you want to measure on your own agent. Benchmarking compares; evaluation answers. Benchmarking is shared; evaluation is yours. This chapter draws the distinction in four ways --- purpose, audience, data, cadence --- and sets up the three chapters that follow: how to design an evaluation framework, how to measure what matters in production, and how to evaluate the things that fail rarely and badly.]{Evaluation vs.\ Benchmarking}
\label{ch:40-evaluation-13-vs-benchmarking}

\starttwocol

\section{The two words your team has been using interchangeably}

Maya's team has been using ``benchmark'' and ``evaluation'' as synonyms for a year. So has every vendor she has met. So has every analyst report on her desk. The reason the distinction matters now, and not earlier, is that her auditors are about to start using the words differently --- and her team's reports will be graded on whether they answer the right question.

A benchmark is a shared test. Somebody else built it, somebody else ran it, and the score it produces is meaningful only against other scores on the same test. SWE-bench Verified \citep{jimenez2024swebench, openai-swebench-verified-2024} is a benchmark. So is $\tau$-bench \citep{yao2024taubench}. So is BFCL \citep{bfcl-2024}. The point of running a benchmark is to place your agent on a leaderboard, or to compare your candidate vendor against the one you already have. The number is comparative or it is nothing.

An evaluation is a purpose-built test of \textit{your} agent doing \textit{your} job. You designed the inputs. You decided what counts as a pass. You run it on every release and you keep the results in the same binder where your audit committee will read them. The number is decision-relevant or it is theatre.

\begin{TNdefinition}{Evaluation}
Pillar 4 of ARIA. A purpose-built, organization-specific test of an agent doing the job you intend it to do. An evaluation answers the question ``is this agent ready for the work we are about to give it?'' A benchmark, by contrast, answers ``how does this agent compare to others on a shared test?''
\end{TNdefinition}

The most common failure pattern in early agent programs is to confuse one for the other. A team picks an agent on the back of a benchmark score, deploys it, and then discovers in the first month of production that the benchmark measured something adjacent to --- but not the same as --- the job at hand. The other failure pattern is to build a sprawling internal evaluation suite that scores well on everything and predicts nothing, because nobody decided what a score of 92\% was supposed to mean.

\begin{TNinpractice}{Northwind Bank}
\textit{Northwind's auditors asked for ``the benchmark report'' on the wealth-management assistant the bank was certifying.} The team sent over a $\tau$-bench score with a polite note explaining the methodology. The audit lead read it, asked a single question --- \textit{does this benchmark measure the agent's performance on Northwind's customers?} --- and waited. It did not. $\tau$-bench is generic. Within two weeks the team had built its first purpose-built evaluation: 240 cases drawn from real call-center transcripts, with sign-off from compliance on what counted as a correct refusal. The benchmark was still useful; it told Northwind how its agent compared to the field. The evaluation was what the auditors actually wanted, because it told Northwind whether the agent could do \textit{this} job for \textit{these} customers.
\end{TNinpractice}

\section{Four axes that distinguish them}

The cleanest way to keep benchmark and evaluation separate in a meeting is to think along four axes. They map onto questions Maya will hear at her review board.

\textbf{Purpose.} A benchmark compares. An evaluation answers. The benchmark score lets you say ``this agent is in the top quartile of code-fixing agents.'' The evaluation lets you say ``this agent resolves 82\% of the support tickets we paid it to resolve, with a 4\% escalation rate and zero refund errors above \$500.'' One of those sentences belongs in a vendor's pitch deck. The other belongs in an audit binder.

\textbf{Audience.} A benchmark is read by the field. An evaluation is read by your audit committee, your CISO, your regulator, and the team on call when the agent misbehaves. The audience for a benchmark is everyone interested in agents. The audience for an evaluation is the small group of people who have to defend a deployment decision.

\textbf{Data.} A benchmark uses public test sets --- the same ones every other vendor runs against. An evaluation uses a \textit{golden dataset} you built yourself, drawn from your own traffic, labelled by your own subject-matter experts, and protected from your model's training data. The benchmark's data is shared and stale; yours is private and fresh.

\textbf{Cadence.} A benchmark gets a run on every model release, every vendor swap, every quarter that you want to recheck the market. An evaluation runs continuously: on every commit, on every prompt change, on every model upgrade, on a sample of live traffic in production. The benchmark cadence is event-driven; the evaluation cadence is operational.

\begin{Figure}
\renewcommand{\arraystretch}{1.2}
\setlength{\tabcolsep}{4pt}
{\footnotesize
\begin{tabularx}{\columnwidth}{@{}>{\bfseries\color{TNNavy}}l X X@{}}
\toprule
 & \textcolor{TNTealDeep}{\bfseries Benchmark} & \textcolor{TNTealDeep}{\bfseries Evaluation} \\
\midrule
Purpose  & Compare models across the field & Defend a deployment decision \\
Audience & Researchers, vendors             & CISO, audit, regulator \\
Data     & Public, shared test sets         & Private golden set from your traffic \\
Cadence  & Per model release                 & Continuous: per commit, per prompt \\
\bottomrule
\end{tabularx}}
\captionof{figure}{Benchmark vs.\ evaluation across the four axes that matter when a vendor's slide arrives on your desk.}
\end{Figure}

\begin{TNwarning}
The vendor's deck is engineered to make a benchmark look like an evaluation. The slide will name the benchmark, quote the score, and add a sentence like ``demonstrates production readiness.'' A leaderboard does not demonstrate production readiness on your traffic. Two questions catch this every time: \textit{did this test use our data?} and \textit{would my audit committee accept this number?} If the answer to either is ``no,'' the slide is a benchmark dressed as an evaluation, and the conversation needs to move.
\end{TNwarning}

\section{Why the distinction goes missing}

There are two structural reasons benchmark and evaluation get confused. The first is that benchmarks are easier to talk about. A leaderboard fits on a slide. A purpose-built evaluation is a small team doing unglamorous work on data your CISO will not let you share externally. When a vendor wants to sell you an agent, they show you the slide.

The second reason is that until 2023 or so, the two largely overlapped. The state of the field was such that if a model topped a public benchmark, it would probably work for your job --- because the public benchmarks were the only meaningful tests, and the jobs were close enough to what the benchmarks measured. By 2024 that stopped being true. Agents started succeeding on benchmarks \citep{jimenez2024swebench, mialon2023gaia, yao2024taubench} and failing in production for reasons benchmarks could not catch. The failure modes were specific: wrong-format tool calls in your specific stack, retrieval that worked on Wikipedia and not on your wiki, polite users in the benchmark and frustrated ones in your call centre.

Two recent surveys make the gap concrete. \citet{mohammadi2025survey} and \citet{yehudai2025survey} both find that benchmark performance correlates weakly with production performance for tool-using agents, and that the gap is widest exactly where Maya cares most: multi-turn conversations, error recovery, and consistency across repeated runs. The CLEAR framework \citep{mehta2025clear} goes further and recommends that any serious enterprise build a purpose-built evaluation suite from day one and treat the public benchmarks as a sanity check, not a decision.

The practical implication for Maya is simple. The benchmark numbers are reference points, not decisions. The decisions live in the evaluation.

\section{The taxonomy that organizes the rest of this part}

The rest of Part IV is structured around three modes of evaluation, mapped onto the agent's lifecycle. The taxonomy below is the one your team should adopt as a default. There are three modes, and most enterprise agents need at least some of each.

\begin{TNdefinition}{Offline evaluation}
Running the agent against a curated, versioned test set before any production traffic touches it. The quality gate. Catches functional regressions, schema violations, and known-bad inputs before they reach users.
\end{TNdefinition}

\begin{TNdefinition}{Online evaluation}
Measuring the agent on a sample of live production traffic, where the inputs are real, the consequences are real, and the ground truth is harder to come by. Catches things the offline test set could not predict: distribution shift, frustrated users, the rare adversarial input.
\end{TNdefinition}

\begin{TNdefinition}{Continuous evaluation}
The closed loop that connects the two. Production traces that fail are harvested into the offline test set. The offline test set, in turn, gets re-run on every model change. Without the loop, both modes go stale.
\end{TNdefinition}

The three modes are not substitutes for each other. Offline tells you whether the agent works on the cases you expect; it cannot tell you what production looks like. Online tells you what production looks like; without offline regression tests, you cannot tell what changed. Continuous keeps both honest. A serious evaluation program runs all three.

The three modes intersect with three things you measure: did the agent do the job (functional, see \autoref{ch:40-evaluation-15-production}), was the answer any good (quality and reliability, also \autoref{ch:40-evaluation-15-production}), and did it stay safe across populations and over time (safety, fairness, and ongoing checks, \autoref{ch:40-evaluation-16-safety}). The taxonomy is a 3-by-3 grid, but no enterprise builds all nine cells at once. \autoref{ch:40-evaluation-14-framework} is about which cells to start with.

\begin{Figure}
\resizebox{\textwidth}{!}{%
\begin{tikzpicture}[
  font=\sffamily\footnotesize,
  cell/.style={draw=TNNavy, line width=0.5pt, rounded corners=2pt,
    fill=TNNavyTint, text width=3.2cm, align=left, inner sep=5pt,
    minimum height=1.85cm, anchor=north west},
  head/.style={font=\sffamily\bfseries\small\color{TNTealDeep}, anchor=center},
  rowlabel/.style={font=\sffamily\bfseries\small\color{TNNavy},
    text width=2cm, align=right, anchor=east}
]
  \node[head] at (1.6, 0.5)  {Offline};
  \node[head] at (5.0, 0.5)  {Online};
  \node[rowlabel] at (-0.1, -1.05) {Functional};
  \node[rowlabel] at (-0.1, -3.05) {Quality \&\\safety};
  \node[cell] at (0,  -0.2) {exact-match,\\schema validity,\\tool-call success};
  \node[cell] at (3.4,-0.2) {task completion,\\time-to-recover,\\escalation rate};
  \node[cell] at (0,  -2.2) {faithfulness,\\groundedness,\\toxicity gate};
  \node[cell] at (3.4,-2.2) {drift, fairness,\\pass\textasciicircum k stability};
\end{tikzpicture}%
}
\captionof{figure}{The evaluation taxonomy. Rows distinguish functional checks from quality and safety; columns separate offline tests from on-traffic measurement. Continuous evaluation is the loop between them.}
\end{Figure}

\section{What this part will and will not give you}

This part is about how to build an evaluation program your auditors, your CISO, and your CFO can read. It is not a survey of every published evaluation paper. It is not a vendor comparison. It is not a list of tools.

You will get the four-axis distinction between benchmark and evaluation, the taxonomy of offline / online / continuous, the structure of a framework that produces evidence rather than vanity numbers, the metrics that matter in production (faithfulness, groundedness, escalation, time-to-recover, \textit{pass\textasciicircum k}), and the safety and fairness measures that turn up in the audit binder. You will get a closing checklist per chapter that gives your team the language to argue about thresholds without arguing about what they mean. What you will not get is a number to put on a slide. The whole point of this part is that the number on the slide is the answer to the wrong question.

The auditor's question --- \textit{how do we know this agent is ready?} --- is the one Part IV teaches you to answer. The answer will not be a single score. It will be a small set of numbers, taken from your own data, on a cadence your team can defend.

\TNrule

\stoptwocol

\begin{TNkeytakeaways}
\item A benchmark compares your agent against others on a shared test; an evaluation answers whether your agent is ready for the work you are about to give it.
\item The four axes that keep them separate are purpose, audience, data, and cadence --- and the audience axis is the one your audit committee cares about most.
\item Benchmarks correlate weakly with production performance for tool-using agents; treat them as reference points, not as decisions.
\item The evaluation taxonomy used through this part has three modes --- offline, online, continuous --- and three concerns --- functional, quality, safety. Your program will need cells from each.
\item The number on the vendor's slide is rarely the number your auditors want. The next three chapters are about producing the number they do.
\end{TNkeytakeaways}

\begin{TNchecklist}
\item Ask your team to label every metric in their current dashboard either ``benchmark'' or ``evaluation,'' and to defend the label in one sentence.
\item For each production agent, write down which questions your purpose-built evaluation must answer --- and the audience for the answer.
\item Confirm that your golden dataset is private, fresh, and drawn from your own traffic (not the vendor's curated examples).
\item Decide which of the three evaluation modes --- offline, online, continuous --- your program runs today, and where the gaps are.
\item In the next vendor meeting, ask the four-axis question: \textit{purpose, audience, data, cadence --- which of these is your benchmark slide actually addressing?}
\item Schedule a 30-minute conversation with audit on what they expect to see in an evaluation report, before the team builds the report.
\end{TNchecklist}



